\documentclass[12pt,a4paper]{article}
\usepackage[utf8]{inputenc}
\usepackage[spanish]{babel}
\usepackage{graphicx}
\usepackage{listings}
\usepackage{xcolor}
\usepackage{hyperref}
\usepackage{minted}

\title{Documentación Técnica - TENA\_FIT}
\author{Angel}
\date{\today}

\begin{document}

\maketitle
\tableofcontents

\section{Introducción}
TENA\_FIT es una aplicación completa de fitness y bienestar que integra tecnologías modernas para proporcionar una experiencia multiplataforma. Este documento técnico detalla la arquitectura, implementación y guías de desarrollo del proyecto.

\section{Arquitectura del Sistema}

\subsection{Frontend (React + Vite)}
\begin{itemize}
    \item \textbf{Framework:} React 18 con TypeScript
    \item \textbf{Bundler:} Vite
    \item \textbf{UI Framework:} Chakra UI
    \item \textbf{Estado:} Context API + Custom Hooks
    \item \textbf{Routing:} React Router DOM v6
    \item \textbf{HTTP Client:} Axios
\end{itemize}

\subsection{Backend (Node.js)}
\begin{itemize}
    \item \textbf{Runtime:} Node.js
    \item \textbf{Framework:} Express
    \item \textbf{Base de Datos:} Sistema de archivos JSON
    \item \textbf{Autenticación:} JWT + bcrypt
    \item \textbf{API:} RESTful
\end{itemize}

\subsection{Aplicación Móvil (Flutter)}
\begin{itemize}
    \item \textbf{Framework:} Flutter SDK
    \item \textbf{Gestión de Estado:} Provider
    \item \textbf{Almacenamiento Local:} SQLite
    \item \textbf{Networking:} HTTP package
\end{itemize}

\section{Base de Datos JSON}

\subsection{Estructura de Archivos}
\begin{itemize}
    \item \texttt{/src/data/users.json} - Información de usuarios
    \item \texttt{/src/data/ejercicios.json} - Catálogo de ejercicios
    \item \texttt{/src/data/rutinas.json} - Rutinas predefinidas
    \item \texttt{/src/data/progreso.json} - Seguimiento de progreso
\end{itemize}

\subsection{Esquemas JSON}

\subsubsection{users.json}
\begin{minted}{json}
{
  "users": [
    {
      "id": "string",
      "nombre": "string",
      "email": "string",
      "password": "string",
      "rol": "admin|cliente"
    }
  ]
}
\end{minted}

\subsubsection{ejercicios.json}
\begin{minted}{json}
{
  "ejercicios": [
    {
      "id": "string",
      "nombre": "string",
      "categoria": "string",
      "descripcion": "string",
      "imagen": "string",
      "musculos": ["string"],
      "nivel": "principiante|intermedio|avanzado"
    }
  ]
}
\end{minted}

\subsubsection{rutinas.json}
\begin{minted}{json}
{
  "rutinas": [
    {
      "id": "string",
      "nombre": "string",
      "descripcion": "string",
      "nivel": "string",
      "duracion": "string",
      "dias_semana": ["string"],
      "ejercicios": [
        {
          "ejercicio_id": "string",
          "series": "number",
          "repeticiones": "number",
          "descanso": "string",
          "notas": "string"
        }
      ]
    }
  ]
}
\end{minted}

\subsubsection{progreso.json}
\begin{minted}{json}
{
  "progreso": [
    {
      "id": "string",
      "usuario_id": "string",
      "fecha": "date",
      "tipo": "peso|ejercicio|objetivo|ejercicio_asignado",
      "valor": "number",
      "unidad": "string"
    }
  ]
}
\end{minted}

\section{Guía de Instalación}

\subsection{Configuración del Entorno de Desarrollo}
\begin{enumerate}
    \item Instalar Node.js (v18.x o superior)
    \item Instalar Flutter SDK (v3.2.3 o superior)
    \item Configurar IDE (VS Code recomendado)
    \item Instalar extensiones necesarias:
        \begin{itemize}
            \item Flutter
            \item Dart
            \item ES7+ React/Redux/React-Native snippets
            \item Thunder Client (para pruebas de API)
        \end{itemize}
\end{enumerate}

\section{API Endpoints}

\subsection{Autenticación}
\begin{itemize}
    \item \texttt{POST /api/auth/login}
    \item \texttt{POST /api/auth/register}
    \item \texttt{POST /api/auth/refresh-token}
\end{itemize}

\subsection{Usuarios}
\begin{itemize}
    \item \texttt{GET /api/users/profile}
    \item \texttt{PUT /api/users/profile}
    \item \texttt{GET /api/users/statistics}
\end{itemize}

\subsection{Ejercicios}
\begin{itemize}
    \item \texttt{GET /api/exercises}
    \item \texttt{POST /api/exercises}
    \item \texttt{PUT /api/exercises/:id}
    \item \texttt{DELETE /api/exercises/:id}
\end{itemize}

\section{Seguridad}

\subsection{Autenticación}
El sistema utiliza JWT (JSON Web Tokens) para la autenticación:
\begin{itemize}
    \item Access Token (15 minutos)
    \item Refresh Token (7 días)
    \item Passwords hasheados con bcrypt
\end{itemize}

\subsection{Variables de Entorno}
\begin{minted}{bash}
# Frontend (.env)
VITE_API_URL=http://localhost:3000/api

# Backend (.env)
PORT=3000
JWT_SECRET=your_jwt_secret
\end{minted}

\section{Credenciales de Administrador}
\begin{minted}{text}
Usuario: admin@tenafit.com
Contraseña: admin123
\end{minted}

\section{Flujo de Trabajo de Desarrollo}

\subsection{Branches}
\begin{itemize}
    \item \texttt{main} - Producción
    \item \texttt{develop} - Desarrollo
    \item \texttt{feature/*} - Nuevas características
    \item \texttt{bugfix/*} - Correcciones
\end{itemize}

\subsection{Commits}
Formato de commits:
\begin{minted}{text}
<tipo>(<alcance>): <descripción>

tipos: feat, fix, docs, style, refactor, test, chore
\end{minted}

\section{Despliegue}

\subsection{Frontend}
\begin{minted}{bash}
cd frontend
npm run build
# Los archivos de distribución estarán en /dist
\end{minted}

\subsection{Backend}
\begin{minted}{bash}
cd backend
npm run build
# Usar PM2 o similar para producción
\end{minted}

\subsection{App Móvil}
\begin{minted}{bash}
cd app_android
flutter build apk --release
# El APK estará en build/app/outputs/flutter-apk/
\end{minted}

\section{Mantenimiento}

\subsection{Logs}
\begin{itemize}
    \item Frontend: Console logs + Sentry
    \item Backend: Winston logger
    \item App Móvil: Flutter logs
\end{itemize}

\subsection{Monitoreo}
\begin{itemize}
    \item Métricas de servidor
    \item Rendimiento de la aplicación
    \item Errores y crashes
\end{itemize}

\section{Conclusión}
Esta documentación proporciona una visión general técnica del proyecto TENA\_FIT. Para más detalles específicos, consultar el código fuente y los comentarios en cada componente.

\end{document} 